\lysh{15993_SOLUTION}{1}
α) Είναι γνωστό ότι η εξίσωση $(x-x_0)^2+(y-y_0)^2=\rho^2$ παριστάνει κύκλο με κέντρο το $K(x_0,y_0)$ και ακτίνα $\rho > 0$.
Η (1) γράφεται
\[ (x-2)^2+(y-\lambda)^2=\left(\sqrt{\lambda^2+1}\right)^2 \]
επομένως παριστάνει κύκλο με κέντρο το $K(2,\lambda)$ και ακτίνα $\rho = \sqrt{\lambda^2+1}$, διότι $\lambda^2+1 > 0$ για κάθε $\lambda \in \mathbb{R}$.

β) Επιλέγουμε δύο από τους κύκλους (1), δίνοντας τις παρακάτω τιμές:
Για $\lambda=0$, $x^2+y^2-4x-3=0$
Για $\lambda=1$, $x^2+y^2-4x-2y+3=0$
Αφαιρώντας κατά μέλη τις δύο ισότητες προκύπτει $y=0$ και αντικαθιστώντας στην πρώτη εξίσωση, είναι: $x^2-4x+3=0 \Leftrightarrow x=1$ ή $x=3$.
Επομένως οι κύκλοι έχουν δύο κοινά σημεία, τα $A(1,0)$ και $B(3,0)$.
Με μια απλή αντικατάσταση στην (1), αποδεικνύεται ότι τα σημεία αυτά την επαληθεύουν για κάθε $\lambda \in \mathbb{R}$ και ως εκ τούτου, αποτελούν τα κοινά σημεία όλων των κύκλων.

γ) Η κοινή χορδή των κύκλων (1) είναι το ευθύγραμμο τμήμα $AB$, το οποίο βρίσκεται πάνω στον άξονα $x'x$. Επομένως έχει εξίσωση $y=0$.
Τα κέντρα όλων των κύκλων είναι της μορφής $K(2,\lambda)$ με $\lambda \in \mathbb{R}$. Άρα, η ευθεία που διέρχεται από τα κέντρα όλων των κύκλων, είναι η κατακόρυφη ευθεία με εξίσωση $x=2$.
Επομένως είναι κάθετη στην κοινή χορδή.

% TODO: TikZ σχήμα

δ) Αφού το σημείο $M(\alpha, \beta)$ επαληθεύει την (1) για κάθε $\lambda \in \mathbb{R}$, πρέπει υποχρεωτικά να είναι ή το $A(1,0)$ ή το $B(3,0)$. Σε κάθε περίπτωση ισχύει: $\alpha \cdot \beta = 1 \cdot 0 = 3 \cdot 0 = 0$.
\endlysh